%-*-latex-*-
%-*-latex-*-
\newcommand\COURSE{ciss350}
\newcommand\ASSESSMENT{q10305}
\newcommand\ASSESSMENTTYPE{Quiz}
\newcommand\POINTS{\textwhite{xxx/xxx}}

\input{myquizpreamble}
\input{yliow}
\input{\COURSE}
%-*-latex-*-

%-*-latex-*-

%-*-latex-*-
\renewcommand\TITLE{\ASSESSMENTTYPE \ \ASSESSMENT}

\renewcommand\EMAIL{}
\renewcommand\AUTHOR{eagraham1@cougars.ccis.edu}

\textwidth=6in
\begin{document}
\topmatterthree

Here's an example on sorting a vector:
\VerbatimInput[frame=single,fontsize=\footnotesize]{sort.cpp}
This sorts the whole vector.
(Note that \verb!std::sort! is not stable.)
For details on \verb!begin()! and \verb!end()!, see
chapter on Containers (on iterators).
\newpage

%------------------------------------------------------------------------------
\nextq
Party planning problem.
There are $n=5$ friends living on a street at address $0, 5, 3, 9, 3$.
The address are integers from $0$ to $N - 1$
where $N$ is huge, say $N = 1000000$.
The houses are on a straight line.
\begin{Verbatim}[frame=single,fontsize=\footnotesize]
            P
            P
P                   P               P
+---+---+---+---+---+---+---+---+---+---+---+---+---+ ...
0   1   2   3   4   5   6   7   8   9   10  11  12  13 ...
\end{Verbatim}
(\verb!P! marks an address where a friend lives there.)
The distance between 0 and 1 is 1, the distance between 1 and 2 is 1, etc.
The friends are planning a party and they have decided to hold it at the house
at address 5.
The total distance travelled by everyone to 5 is $5+2+2+0+4 = 13$.
\begin{Verbatim}[frame=single,fontsize=\footnotesize]
            P------>
            P------>
P------------------>P<--------------P
+---+---+---+---+---+---+---+---+---+---+---+---+---+ ...
0   1   2   3   4   5   6   7   8   9   10  11  12  13 ...
\end{Verbatim}
Find \lq\lq best house", i.e.,
the one where the total distance travelled is the smallest.
Note that the brute force solution (using CISS240 methods) would run in $O(n^2)$.
Your algorithm must run with time of $O(n \lg n)$ and a space complexity
of $O(1)$.
(There's actually a faster algorithm that runs with $O(n)$ time.)
\\
\ANSWER
\begin{answercode}
#include <iostream>
#include <cstdlib>
#include <vector>
#include <algorithm>

int main()
{
    int seed;
    std:cin >> seed;
    srand(seed);
    int n = 1000000;
    std::vector< int > x(n);
    for (int i = 0; i < n; ++i)
    {
        x[i] = rand() % 1000000;
    }
    int best_address = x[0];

    std::sort(x.begin(), x.end());

    int best_address = x[n / 2];

    std::cout << best_address << '\n';
}
\end{answercode}

%------------------------------------------------------------------------------
\nextq
A permutation is a shuffle of values.
For instance 1, 3, 5, 6, 4 is a permutation of 3, 1, 4, 6, 5.
Complete the following problem.
Clearly, your code that sets \verb!is_permutation! to \verb!true!
if and only if \verb!u! is a permutation of \verb!v!.
This must work for different vectors including vectors
of large sizes.
You should not assume the values are unique.
Your algorithm must run with time of $O(n \lg n)$ and a space complexity
of $O(1)$.
\\
\ANSWER
\begin{answercode}
#include <iostream>
#include <cstdlib>
#include <vector>
#include <algorithm>

int main()
{
    std::vector< int > u {1, 3, 5, 6, 4};
    std::vector< int > v {3, 1, 4, 6, 5};
    bool is_permutation = true;

    if(u.size() != v.size())
    {
      is_permutation = false;
    }
    else
    {
      std::sort(u.begin(), u.end());
      std::sort(u.begin(), u.end());

      if(u[i] != v[i])
      {
        is_permutation = false;
        break;
      }
    }

    std::cout << is_permutation << '\n';
}
\end{answercode}

%------------------------------------------------------------------------------
%-*-latex-*-
\newpage
\input{instructions}
\end{document}

